\documentclass[10pt, a4paper]{article}
\usepackage[left=2cm, right=2cm, top=2.5cm, bottom=2.5cm]{geometry}
\usepackage{longtable}
\usepackage{array}
\usepackage{xcolor}
\usepackage{colortbl}
\usepackage{fancyhdr}
\usepackage{booktabs}
\usepackage[hidelinks]{hyperref}
\setlength{\emergencystretch}{3em}

\definecolor{passgreen}{RGB}{34,139,34}
\definecolor{failred}{RGB}{200,0,0}
\definecolor{warnorange}{RGB}{200,100,0}
\definecolor{skipgray}{RGB}{120,120,120}

\title{\textbf{MEA OCPP 1.6 Compliance Test Report}\\[0.3em]
\large Section 9: MEA-Specific Configuration \& V2G/BPT\\[0.3em]
\normalsize Vector vSECC.single Board vs MEA CSMS (rddQC4000001)}
\author{MEA-V2G Project --- Automated Test}
\date{2026-06-09 15:50}

\pagestyle{fancy}
\fancyhf{}
\lhead{MEA OCPP 1.6 -- Section 9}
\rhead{Page \thepage}

\begin{document}
\maketitle
\tableofcontents
\newpage

\section{Test Configuration}
\begin{tabular}{ll}
\textbf{Device Under Test}   & Vector vSECC.single Board \\
\textbf{Device IP}           & 192.168.1.166 \\
\textbf{OCPP Version}        & 1.6 (JSON) \\
\textbf{CP ID}               & rddQC4000001 \\
\textbf{CSMS}                & wss://ocpp.measandbox.com:2930 \\
\textbf{Connection}          & Direct (no proxy) \\
\textbf{Test PC}             & 192.168.111.185 (alias 192.168.1.200/24 on enp3s0) \\
\textbf{MeterValues Timeout} & 15 s per power-demand item \\
\textbf{Test Date}           & 2026-06-09 08:44:14 UTC \\
\textbf{Results file}        & tex/vsecc\_section9\_results.json \\
\end{tabular}

\subsection*{Test Architecture}
The test calls the MEA sandbox REST API with HTTP Digest authentication to issue
\texttt{ChangeConfiguration} commands, then subscribes to the vSECC MQTT broker
to observe \texttt{metervalues} publications as confirmation that power-demand
setpoints were received by the charger.

\subsection*{MEA-Specific OCPP Keys}
The MEA CSMS uses two non-standard OCPP~1.6 configuration keys for V2G/BPT control:
\begin{itemize}
  \item \textbf{V2GMode} --- Enables bidirectional power transfer (BPT) capability
        at the CSMS level. Value: \texttt{"true"} or \texttt{"false"}.
  \item \textbf{Power.Active.Import} --- Sets the instantaneous power demand in Watts.
        Negative values request discharge from EV to grid (V2G).
        Positive values request charge from grid to EV.
        Zero stops power flow.
\end{itemize}
These keys are not defined in the OCPP~1.6 specification.
They are MEA-specific extensions used for V2G/BPT scheduling and control.

\section{Section 9 Results}
\textbf{Summary: 5 PASS \quad 0 FAIL \quad 3 WARN \quad 0 SKIP \quad (8 total)}

\renewcommand{\arraystretch}{1.35}
\newcolumntype{L}[1]{>{\raggedright\arraybackslash}p{#1}}
\newcolumntype{M}[1]{>{\small\ttfamily\raggedright\arraybackslash}p{#1}}
\newcolumntype{C}[1]{>{\centering\arraybackslash}p{#1}}
\begin{longtable}{|L{1.2cm}|L{5.5cm}|M{6.8cm}|C{1.8cm}|}
\hline
\rowcolor{gray!25}
\textbf{Item} & \small\textbf{Test Description} & \small\textbf{Detail / Evidence} & \small\textbf{Result} \\
\hline
\endfirsthead
\hline
\rowcolor{gray!25}
\textbf{Item} & \small\textbf{Test Description} & \small\textbf{Detail / Evidence} & \small\textbf{Result} \\
\hline
\endhead
\hline
\endfoot
\small 9.0 & \small CSMS connection probe — GetConfiguration HTTP 200 & vSECC connected to MEA CSMS & \textcolor{passgreen}{\textbf{PASS}} \\
\hline
\multicolumn{4}{|l|}{\cellcolor{gray!10}\small\textbf{9.1--9.3 Configuration (Standard Keys via ChangeConfiguration)}} \\
\hline
\small 9.1 & \small ChangeConfiguration V2GMode=true → Accepted & HTTP 200 & \textcolor{passgreen}{\textbf{PASS}} \\
\hline
\small 9.2 & \small ChangeConfiguration MeterValueSampleInterval=10 → Accepted & HTTP 200 & \textcolor{passgreen}{\textbf{PASS}} \\
\hline
\small 9.3 & \small ChangeConfiguration LocalAuthorizeOffline=false → Accepted & HTTP 200 & \textcolor{passgreen}{\textbf{PASS}} \\
\hline
\multicolumn{4}{|l|}{\cellcolor{gray!10}\small\textbf{9.4 Power Demand (V2G/BPT) --- MEA-Specific Keys}} \\
\hline
\small 9.4.0 & \small ChangeConfiguration V2GMode=true (pre-condition for power demand) & HTTP 200 & \textcolor{passgreen}{\textbf{PASS}} \\
\hline
\small 9.4.1 & \small Power Demand -5000 W (BPT discharge from EV to grid) & HTTP 200; metervalues not observed in 15s{\newline\normalfont\footnotesize\itshape V2G/BPT requires active charging session and V2G-capable EV} & \textcolor{warnorange}{\textbf{WARN}} \\
\hline
\small 9.4.2 & \small Power Demand +2000 W (BPT charge EV from grid) & HTTP 200; metervalues not observed in 15s{\newline\normalfont\footnotesize\itshape V2G/BPT requires active charging session and V2G-capable EV} & \textcolor{warnorange}{\textbf{WARN}} \\
\hline
\small 9.4.3 & \small Power Demand 0 W (BPT idle / stop power flow) & HTTP 200; metervalues not observed in 15s{\newline\normalfont\footnotesize\itshape V2G/BPT requires active charging session and V2G-capable EV} & \textcolor{warnorange}{\textbf{WARN}} \\
\hline

\end{longtable}
\normalsize

\section{Notes on Failures and Warnings}

\subsection*{9.1 -- 9.3 Configuration items}
Items 9.1--9.3 call \texttt{POST /EV/remote/changeConfiguration} with standard
OCPP~1.6 keys (\texttt{V2GMode}, \texttt{MeterValueSampleInterval},
\texttt{LocalAuthorizeOffline}). A result of HTTP~200 from the MEA CSMS confirms
the command was forwarded to the charger. If the charger returns
\texttt{Rejected} in the response body, the detail field will show
\texttt{status=Rejected}; this indicates the vSECC does not support that key.

\subsection*{9.4 Power Demand (V2G/BPT) items}
\begin{itemize}
  \item \textbf{9.4.0 -- 9.4.3} use the MEA-specific key
        \texttt{Power.Active.Import} to set an instantaneous power demand.
        The key is not defined in the OCPP~1.6 specification; it is a
        MEA extension for BPT (bidirectional power transfer) scheduling.
  \item \textbf{PASS} is assigned when the MEA API returns HTTP~200,
        indicating the CSMS accepted and forwarded the command.
  \item \textbf{WARN} is assigned when \texttt{metervalues} events are not
        observed on the vSECC MQTT broker within the timeout. This is expected
        when no active charging session exists or when the connected EV does not
        support V2G/ISO~15118~-2 power scheduling.
  \item V2G power control requires: (a)~an active charging session,
        (b)~a V2G-capable EV, (c)~a signed \texttt{PowerDeliveryReq} exchange
        at the ISO~15118 level. The MEA REST API call alone is not sufficient
        to verify end-to-end power delivery.
\end{itemize}

\section{Dependency Map}
\begin{tabular}{ll}
\textbf{Configuration} & 9.0 (connection) $\rightarrow$ 9.1, 9.2, 9.3 (independent) \\
\textbf{V2G/BPT}       & 9.4.0 (V2GMode=true) $\rightarrow$ 9.4.1 $\rightarrow$ 9.4.2 $\rightarrow$ 9.4.3 \\
\end{tabular}

\end{document}
